% Part: first-order-logic 
% Chapter: axiomatic-deduction 
% Section: deduction-theorem-quantifiers

\documentclass[../../../include/open-logic-section]{subfiles}

\begin{document}

\olfileid{fol}{axd}{ddq}

\olsection{De deductiestelling met kwantoren (‘voor alle’ en ‘er bestaat’)}

\begin{thm}[Deductiestelling: zet de extra aanname in de implicatie]
\ollabel{thm:deduction-thm-q} Als $\Gamma \cup \{!A\} \Proves !B$,
dan $\Gamma \Proves !A \lif !B$.
\end{thm}

\begin{proof}
We gebruiken weer inductie naar de lengte van de !!{derivation} van
$!B$ uit $\Gamma \cup \{!A\}$: eerst een afleiding van lengte één,
daarna de stap naar een langere afleiding.

Voor de inductiebasis geldt precies hetzelfde bewijs als in
\olref[ded]{thm:deduction-thm}.

Nu de inductiestap. Neem weer aan dat de !!{derivation} van $!B$ uit
$\Gamma \cup \{!A\}$ eindigt met een stap~$!B$ die met een
afleidingsregel is verkregen. Is die regel modus ponens, dan gebruiken
we hetzelfde bewijs als in \olref[ded]{thm:deduction-thm}. Is die regel
\QR, dan geldt $!B \ident !C \lif \lforall[x][!D(x)]$. Eerder in de
!!{derivation} staat dan !!a{formula} van de vorm $!C \lif !D(a)$.
Daarbij komt $a$ niet voor in~$!C$, niet in~$!A$ en ook niet in
$\Gamma$. We hebben dus:
\begin{align*}
  \Gamma \cup \{!A\} & \Proves !C \lif !D(a),\\
  \intertext{De bijbehorende afleiding is korter. Volgens de inductiehypothese geldt daarom}
    \Gamma & \Proves !A \lif (!C \lif !D(a)).\\
  \intertext{Gebruik nu}
  & \Proves (!A \lif (!C \lif !D(a))) \lif ((!A \land !C) \lif !D(a))\\
  \intertext{Pas modus ponens toe. Dan krijgen we}
  \Gamma & \Proves (!A \land !C) \lif !D(a).\\
  \intertext{De eigenvariabelevoorwaarde geldt nog steeds: $a$ komt nergens voor waar dat verboden is. We mogen dus met \QR{} een stap aan deze !!{derivation} toevoegen. Dat geeft}
    \Gamma & \Proves (!A \land !C) \lif \lforall[x][!D(x)].\\
    \intertext{We hebben bovendien}
    & \Proves ((!A \land !C) \lif \lforall[x][!D(x)]) \lif (!A \lif (!C \lif \lforall[x][!D(x)]),\\
    \intertext{Pas nogmaals modus ponens toe. Dan}
    \Gamma & \Proves !A \lif (!C \lif \lforall[x][!D(x)]),
\end{align*}
dus $\Gamma \Proves !B$.

Er blijft één geval over als oefening: $!B$ is met de regel \QR{}
verkregen en heeft de vorm $\lexists[x][!D(x)] \lif !C$.
\end{proof}

\begin{prob}
  Maak het bewijs van \olref[fol][axd][ddq]{thm:deduction-thm-q} af.
\end{prob}

\end{document}
